\chapter{\label{ap:formulas}APÊNDICE A — Formulação matemática detalhada dos algoritmos}

Este apêndice reúne as formulações matemáticas que descrevem o funcionamento
interno dos métodos de detecção e de correção utilizados no trabalho. Elas foram
deslocadas do corpo do texto por não serem essenciais à compreensão das decisões
de projeto --- que se apoiam nos \emph{valores} medidos (precisão, estabilidade,
latência) e não na dedução das equações ---, ficando aqui disponíveis para
consulta de quem desejar reproduzir os cálculos ou inspecionar os detalhes de
implementação.

\section{\label{ap:autocorr_yin}Autocorrelação e função de diferença do YIN}

A autocorrelação discreta de um quadro do sinal (Seção~\ref{sec:autocorrelacao}),
usada para localizar o atraso correspondente ao período fundamental, é dada por:

\begin{equation}
r_t(\tau) = \sum_{j=t+1}^{t+W} x_j\, x_{j+\tau},
\end{equation}

\noindent em que $x_j$ representa o sinal de entrada, $W$ é o tamanho da janela de
análise e $\tau$ é o atraso considerado. O YIN (Seção~\ref{sec:yin}) substitui a
busca por máximos de similaridade por uma função de diferença, que mede a
dissimilaridade entre o sinal e sua versão deslocada:

\begin{equation}
d_t(\tau) = \sum_{j=1}^{W} \bigl(x_j - x_{j+\tau}\bigr)^2 .
\end{equation}

\section{\label{ap:cmndf}CMNDF e candidatos do pYIN}

A função de diferença cumulativa normalizada (CMNDF), empregada na detecção de
pitch do protótipo (Seção~\ref{sec:impl_dsp}), normaliza a função de diferença
para suprimir mínimos espúrios em atrasos pequenos:

\begin{equation}
d'_t(\tau) = \begin{cases}
1, & \tau = 0 \\[4pt]
\dfrac{d_t(\tau) \cdot \tau}{\displaystyle\sum_{j=1}^{\tau} d_t(j)}, & \tau \geq 1
\end{cases}
\quad \text{onde} \quad
d_t(\tau) = \sum_{j=0}^{W-1} \bigl(x_{t+j} - x_{t+j+\tau}\bigr)^2 .
\end{equation}

O pYIN gera $K = 100$ candidatos de pitch por quadro, um para cada limiar $s_k$ de
uma distribuição. O peso atribuído a cada limiar segue uma distribuição Beta
adaptada, que favorece valores intermediários \cite{mauch2014pyin}:

\begin{equation}
w_k \propto s_k \cdot (1 - s_k)^{17}, \quad k = 1, \ldots, K .
\end{equation}

\section{\label{ap:viterbi}Rastreamento temporal por Viterbi causal}

O passo de Viterbi sobre o modelo oculto de Markov (Seção~\ref{sec:impl_streaming})
atualiza, a cada quadro $q$, a coluna de log-probabilidades acumuladas:

\begin{equation}
\delta_q(s) = \max_{s'} \bigl[\delta_{q-1}(s') + \log p(s' \to s)\bigr]
            + \log p_{\text{obs}}(s \mid \text{obs}_q).
\end{equation}

A estimativa de $F_0$ de cada quadro é emitida retrotraçando $L$ passos pela cadeia
de ponteiros de retrocesso $\psi$, a partir do melhor estado da coluna atual:

\begin{equation}
s_{q-L} = \psi_{q-L+1}\bigl[\psi_{q-L+2}[\ldots\psi_q[s^*]\ldots]\bigr],
\quad s^* = \arg\max_s \,\delta_q(s).
\end{equation}

\section{\label{ap:psola}Síntese TD-PSOLA}

Na síntese TD-PSOLA (Seção~\ref{sec:impl_dsp}), as marcas de análise são posicionadas
uma por período; a partir da primeira (ancorada no pico de maior amplitude), cada
marca seguinte é localizada por correlação cruzada com o entorno da anterior,
garantindo coerência de fase entre grãos:

\begin{equation}
m_{k+1} = \arg\max_{c \,\in\, [m_k + T - T/4,\; m_k + T + T/4]}
          \sum_{j=-T/2}^{T/2} x[m_k + j] \cdot x[c + j],
\end{equation}

\noindent onde $T = \lfloor f_s / F_0[m_k] \rceil$ é o período estimado em amostras.
Os grãos de síntese são então sobrepostos por uma janela de Hann de meia-largura
$L = T_{\text{ana}}$:

\begin{equation}
y[\hat{m} + k] \mathrel{+}= x[m + k] \cdot w(k), \quad
w(k) = \tfrac{1}{2}\!\left(1 - \cos\!\frac{\pi(k+L)}{L}\right),
\end{equation}

\noindent em que $m$ é a marca de análise e $\hat{m}$ a posição de síntese
correspondente.

\section{\label{ap:guarda}Folga de guarda da síntese incremental}

A fronteira de decisão da síntese incremental (Seção~\ref{sec:impl_streaming})
desconta uma folga de guarda de dois períodos do tom mais grave configurado, o que
garante que o último grão de cada janela nunca alcance a região ainda não
finalizada:

\begin{equation}
p_{\text{guard}} = \left\lfloor \frac{2\, f_s}{F_{\min}} \right\rfloor .
\end{equation}

\section{\label{ap:alinhamento}Alinhamento das trajetórias no estudo comparativo}

No estudo comparativo (Seção~\ref{sec:pipeline_experimental}), como cada algoritmo
entrega a F$_0$ em sua própria grade temporal, a referência do \textit{Vocadito} é
reamostrada por interpolação linear sobre a grade do algoritmo antes do cálculo das
métricas:

\begin{equation}
\label{eq:alinhamento}
f_{\text{ref}}[i] \;=\; \mathrm{interp}\!\left(t_i;\; t_{\text{src}},\; f_{\text{src}}\right),
\end{equation}

\noindent em que $t_{\text{src}}$ e $f_{\text{src}}$ são os instantes e valores
originais da anotação.
